\section{Decomposition and exhaustion}\label{sec:decomposition-exhaustion}

% Mathematical job:
%   - Define the decomposition record Dec(D).
%   - State six-channel exactness: Dec(D) has exactly six role channels
%     P1 through P6, each with exactly one status from
%       {active_projection, below_threshold, collapsed_boundary_case,
%        absent_with_record, inapplicable_outside_scope}.
%   - State the channel-vs-activation discipline explicitly: inactive
%     statuses are statuses, NOT active witnesses.
%   - State residual exhaustion: every raw feature not absorbed into an
%     active projection or status record is classified by the repaired
%     upper-bound categories.
%   - Prove existence of Dec(D) for every D in FATCD.
%   - Make clear: decomposition is EXISTENTIAL, not unique.
%
% Governing sources:
%   workspace/meta/lens_lab/campaigns/paper3_exact_six_exhaustiveness/decomposition_exhaustion/exact_six_channel_vs_activation_note.md
%   workspace/meta/lens_lab/campaigns/paper3_exact_six_exhaustiveness/decomposition_exhaustion/decomposition_exhaustion_theorem_spec.md
%   workspace/meta/lens_lab/aristotle/results/paper3/paper3_decomposition_exhaustion_exact_G/paper3_decomposition_exhaustion_exact_G_result_shape.md
%   workspace/meta/lens_lab/papers/paper3/drafting_support/theorem_statement_bank.md
%
% Overclaim risks:
%   - Existential-to-unique decomposition slippage ("the decomposition",
%     "canonical decomposition", "the unique decomposition interface").
%   - Conflating six channels with six active witnesses.
%   - Treating a below-threshold status record as a witness for P_i.


The lower-bound theorem of Section~\ref{sec:lower-bounds} and the
upper-bound classification of
Section~\ref{sec:upper-bound-classification} reduce the scoped
exact-six question to a decomposition-existence problem: given
\(D \in \FATCD{}\), does \(D\) admit a well-formed record combining
the six role channels, their statuses, and the classification of any
raw features not absorbed into active projections? This section
answers affirmatively and records the channel-versus-activation
discipline that the answer requires. The lower bounds show that each
of the six channels is materially realizable in the domain; the
upper-bound classification shows that no covered seventh role channel
is needed.

\subsection{Decomposition records and channel statuses}\label{subsec:dec-records}

\begin{definition}[Channel status]\label{def:channel-status}
For each \(i \in \{1,\ldots,6\}\), the \emph{role channel}
\(\Chan{i}{D}\) of a description \(D \in \FATCD{}\) is the
\(P_i\)-labeled slot in the six-channel interface associated to \(D\),
regardless of whether that slot is active. A \emph{channel status} for
\(\Chan{i}{D}\) is exactly one of the five values
\begin{itemize}[topsep=2pt,itemsep=1pt]
  \item \(\activeproj\) --- an active derived role projection
    \(\pi_i(D)\) in the sense of
    Definition~\ref{def:role-projection} exists, with the full
    attachment data of
    Definition~\ref{def:threshold-attachment};
  \item \(\belowthr\) --- candidate raw features are present but the
    threshold data required by
    Definition~\ref{def:threshold-attachment} is absent;
  \item \(\collapsedbc\) --- the candidate cluster is a boundary case
    that collapses under the witness schema \(W_i\);
  \item \(\absentrec\) --- no relevant raw features for role \(P_i\) are
    present in \(\Raw{D}\), and the absence is itself recorded;
  \item \(\outsidescope\) --- the role-\(P_i\) content of \(D\) lies
    outside the covered scope.
\end{itemize}
We write \(\Status{i}{D}\) for the channel status of \(\Chan{i}{D}\).
\end{definition}

\begin{definition}[Decomposition/exhaustion record]\label{def:decomposition-record}
A \emph{decomposition/exhaustion record} for \(D \in \FATCD{}\) is a tuple
\[
\Dec{D} = \bigl(\Raw{D},\; (\Chan{i}{D}, \Status{i}{D})_{i=1}^{6},\;
\Residual{D},\; \mathrm{Classify},\; \mathcal{L},\; \mathcal{V},\;
\Audit{D},\; \mathcal{N}\bigr)
\]
where
\begin{itemize}[topsep=2pt,itemsep=1pt]
  \item the six role channels together with their statuses form the
    six-channel status sequence. If \(\Status{i}{D} = \activeproj\), the channel
    carries an active projection \(\pi_i(D)\) with its full attachment
    data. If \(\Status{i}{D}\) is inactive, the channel carries the
    corresponding below-threshold, collapsed-boundary,
    absent-with-record, or outside-scope status record;
  \item \(\Residual{D} \subseteq \Raw{D}\) is the residual feature set
    (the raw records and annotations not already accounted for by an
    active projection or an inactive status record), and
    \(\mathrm{Classify}\) assigns to each \(r \in \Residual{D}\) a label
    in the scheme of Definition~\ref{def:residual-scheme};
  \item \(\mathcal{L}\) collects the loss and lift records for any level
    crossings involved in projections or classifications;
  \item \(\mathcal{V}\) collects instrument-visibility and
    empirical-bridge annotations where these are asserted;
  \item \(\Audit{D}\) is the honest bookkeeping subrecord;
  \item \(\mathcal{N}\) is the explicit nonclaim record.
\end{itemize}
A decomposition/exhaustion record is \emph{well-formed} when each
\(\Status{i}{D}\) is drawn from the five values of
Definition~\ref{def:channel-status}, each active
\(\pi_i(D)\) is admissible in the sense of
Definitions~\ref{def:role-projection} and~\ref{def:threshold-attachment},
each inactive status record is honestly supported by the corresponding
raw-feature situation, and every raw feature is accounted for at the
channel/exhaustion level either by an active projection, by a status
record for its channel, or by a residual classification. The records in
\(\mathcal{L}\) and \(\mathcal{V}\) are auxiliary annotations on this
coverage, not additional role channels.
\end{definition}


\subsection{Existence, six-channel exactness, and residual exhaustion}\label{subsec:dec-theorems}

\begin{theorem}[Decomposition existence]\label{thm:decomposition-existence}
Every description \(D \in \FATCD{}\) admits a well-formed decomposition/exhaustion
record \(\Dec{D}\) in the sense of
Definition~\ref{def:decomposition-record}.
\end{theorem}

\begin{proof}
For each \(i \in \{1,\ldots,6\}\), inspect \(\Raw{D}\) for a cluster
aligned with the meaning of \(P_i\) under
Proposition~\ref{prop:alignment-rules}. If such a cluster carries the
threshold, witness, level, and audit data of
Definition~\ref{def:threshold-attachment}, record
\(\Status{i}{D} = \activeproj\) and an active projection
\(\pi_i(D)\) in the sense of Definition~\ref{def:role-projection};
otherwise record one of the four inactive statuses, determined by
which attachment data is missing, whether the cluster is a
boundary-case collapse, whether no relevant cluster exists, or
whether the role lies outside the covered scope. Six status entries
result, one per role. The raw features absorbed by these channel
records are removed from the residual stage; the remainder pass to
\(\Residual{D}\). By Theorem~\ref{thm:upper-bound-classification},
every covered residual feature classifies under categories (ii)--(ix)
of Definition~\ref{def:residual-scheme}, with categories (x) and (xi)
empty in the covered scope, so \(\mathrm{Classify}\) is total on
\(\Residual{D}\). The records \(\mathcal{L}\) and \(\mathcal{V}\) are
read from \(\Raw{D}\) and \(\Audit{D}\); the bookkeeping subrecord is
\(\Audit{D}\); and \(\mathcal{N}\) records the nonclaims carried
forward from the witness discipline. The resulting tuple is a
decomposition/exhaustion record, well-formed by construction.
Non-vacuity of the six channels follows from
Theorem~\ref{thm:six-role-lower-bounds}, which supplies, for each
channel, a description in which that channel is active.
\end{proof}

\begin{theorem}[Six-channel exactness and residual exhaustion]\label{thm:six-channel-exactness}
For every \(D \in \FATCD{}\) and every well-formed
\(\Dec{D}\) in the sense of
Definition~\ref{def:decomposition-record}:
\begin{enumerate}[label=(\roman*),topsep=2pt,itemsep=1pt]
  \item the channel sequence \((\Chan{i}{D})_{i=1}^{6}\) has length
    exactly six, one per primitive role \Pone{}, \ldots, \Psix{};
  \item every raw feature of \(D\) is accounted for at the
    channel/exhaustion level by an active projection \(\pi_i(D)\), a
    channel-status record for an inactive channel, or a residual
    classification under the scheme of
    Definition~\ref{def:residual-scheme}; the records in \(\mathcal{L}\)
    and \(\mathcal{V}\) are auxiliary annotations on this coverage;
  \item within the covered scope, no raw feature is classified as
    an unresolved seventh-role threat or as a genuine seventh
    irreducible typed closure role
    (Theorem~\ref{thm:upper-bound-classification}).
\end{enumerate}
\end{theorem}

\begin{proof}
Exactness of the channel sequence is built into the definition of
\(\Dec{D}\); there are six role channels by
Definition~\ref{def:decomposition-record}. Coverage of raw features is
the well-formedness condition of
Definition~\ref{def:decomposition-record} together with totality of
\(\mathrm{Classify}\) on \(\Residual{D}\), which is delivered by
Theorem~\ref{thm:upper-bound-classification}. Emptiness of categories
(x) and (xi) within the covered scope is
Theorem~\ref{thm:upper-bound-classification} directly.
\end{proof}

\subsection{Channel-versus-activation discipline}\label{subsec:channel-vs-activation}

Theorem~\ref{thm:six-channel-exactness} asserts six channels in every
decomposition/exhaustion record. It does not assert six active
projections in every record, nor does it assert that a well-formed
record is unique. The distinction between channel presence and
channel activation prevents six-channel exactness from being read as
an all-six-active claim.

\begin{proposition}[Channel-versus-activation discipline]\label{prop:channel-vs-activation}
Let \(D \in \FATCD{}\) and let \(\Dec{D}\) be a well-formed decomposition/exhaustion
record. The following are distinct statements:
\begin{enumerate}[label=(\roman*),topsep=2pt,itemsep=1pt]
  \item \emph{Six channels}: the channel sequence
    \((\Chan{i}{D})_{i=1}^{6}\) has length six with the fixed role names
    \Pone{}, \ldots, \Psix{}.
  \item \emph{Six active projections}: for each \(i\),
    \(\Status{i}{D} = \activeproj\) and the channel carries an active
    \(\pi_i(D)\).
\end{enumerate}
Theorem~\ref{thm:six-channel-exactness} establishes (i) for every
\(D \in \FATCD{}\). It does not establish (ii). In particular, an
admissible \(D \in \FATCD{}\) may carry any number of active projections
between zero and six; inactive channels are recorded with status
\(\belowthr\), \(\collapsedbc\), \(\absentrec\), or \(\outsidescope\)
and are not active witnesses for the corresponding roles.
\end{proposition}

\begin{proposition}[Decomposition records are not unique]\label{prop:non-uniqueness}
A description \(D \in \FATCD{}\) may admit more than one well-formed
decomposition/exhaustion record in the sense of
Definition~\ref{def:decomposition-record}. Distinct well-formed records
may differ in level assignments, choice of witness schema, residual
classification labels for features admitting more than one admissible
category under the alignment rules, or placement of loss/lift and
visibility records. Theorems~\ref{thm:decomposition-existence}
and~\ref{thm:six-channel-exactness} assert existence, not uniqueness.
\end{proposition}

\begin{remark}[What the decomposition theorems do and do not prove]\label{rem:dec-scope}
Theorems~\ref{thm:decomposition-existence}
and~\ref{thm:six-channel-exactness} establish the structural side of
the scoped exact-six claim: every \(D \in \FATCD{}\) admits a
six-channel decomposition/exhaustion record with residual exhaustion
in the covered scope. Together with the lower-bound theorem
(Section~\ref{sec:lower-bounds}) and the upper-bound classification
(Section~\ref{sec:upper-bound-classification}), they prepare the
scoped exact-six theorem but do not replace the integrated-stress
check that follows. They do not assert that every description
activates all six roles, that the decomposition is canonical or
unique, or that the scoped exact-six program extends beyond \FATCD{};
those remain nonclaims.
\end{remark}

